\subsection{Deviation plan computation}
\label{subsec:deviation_plan_maker}

After demand generation, the next step is to define how the network can bypass
the closure. The \texttt{deviation\_plan\_maker} takes the SUMO network and a
user-selected closure as input, then returns candidate detour paths around the
disrupted segment. These plans are not evaluated at this stage, they are passed
unchanged to the Monte Carlo simulation.

This module contains the main detour-construction contribution of the thesis:
it converts a closure selected on the SUMO network into explicit detour-plan
objects that can later be spliced into route files. The graph-search routines
inside this module are standard algorithms. Dijkstra's algorithm is used for
shortest paths, Yen's algorithm is used to enumerate loopless alternatives and
breadth-first search (BFS) is used to discover upstream origins. The thesis
contribution is how these algorithms are combined with closure-endpoint
promotion, blocked-edge handling and depth-aware route rewriting.

Formally, let $G = (V, E, \ell)$ be a weighted directed graph in which $V$
is the set of junctions, $E$ the set of road edges and
$\ell : E \to \mathbb{R}_{>0}$ an edge length function. Given a set of
closed edges $E_c \subset E$ and a blocked set $B \supseteq E_c$, a
\emph{deviation plan} is defined as a tuple
\begin{equation}
  \mathcal{D} \;=\; \bigl(s,\;\, t,\;\, P,\;\, \ell(P),\;\, d\bigr),
  \label{eq:deviation_plan}
\end{equation}
where $s \in V$ is the \emph{origin} of this particular plan (the junction
at which the detour begins) and $t \in V$ is the detour target, while
$P = (e_1, \ldots, e_m)$ is a loopless path in $G' = (V,\, E \setminus B)$
from $s$ to $t$ and $\ell(P) = \sum_{i=1}^{m} \ell(e_i)$ is its total
geometric length in meters. The integer $d \in \mathbb{N}_0$ is the
\emph{breadth-first search (BFS) depth} of $s$ relative to the
\emph{effective detour source}
returned by the endpoint computation of
Section~\ref{sssec:detour_endpoints}. In the common case this effective
source coincides with the upstream junction of the first closed edge
$s_0 = \mathit{from}(e_1^c)$. If that junction has no usable alternative exit,
the endpoint computation moves the source further upstream and adds the
traversed prefix edges to the blocked set~$B$. In the upstream method, the
origin $s$ can differ from one plan to another: the effective source, $t$ and
$B$ are shared, but individual plans may start further upstream. When
$d = 0$, the plan starts at the effective source and acts as the shallowest
fallback. Plans with $d > 0$ start further upstream and are tried first during
multi-plan route rewriting.

\subsubsection{Junction-level graph representation}

The \texttt{deviation\_plan\_maker} operates on a \emph{junction-level}
directed graph, in contrast to the edge-expanded representation used by
the scenario generator in Section~\ref{sssec:route_library}. Here
nodes correspond to SUMO junctions and arcs correspond to road edges:
an arc $(u, v)$ exists for every edge whose \texttt{from} attribute is
junction $u$ and whose \texttt{to} attribute is junction $v$. Arc weights
default to the geometric length of the edge in meters. An alternative
weight can be loaded from a SUMO \texttt{edgedata.xml} file, in which
case values are averaged across all measurement intervals and the
resulting mean is stored as the weight, if the field encodes speed
rather than travel time, the reciprocal is stored instead so that
shorter weighted distances correspond to faster edges.

The graph is built directly from the \texttt{.net.xml} file. The implementation keeps
three dictionaries (\texttt{edges}, \texttt{outgoing} and \texttt{incoming})
that provide $O(1)$ lookup of any edge record and $O(\deg(v))$ enumeration of a
node's neighbours\footnote{Here, $\deg(v)$ denotes the number of outgoing
or incoming edges adjacent to node $v$, depending on the direction of the
enumeration. Thus, neighbour enumeration is linear in the local branching
factor of the queried junction, not in the total number of edges in the
network.}. These structures are sufficient for the routing operations used in
this module.

\subsubsection{Detour endpoint computation}
\label{sssec:detour_endpoints}

Given a sequence of closed edges $E_c = (e_1, \ldots, e_n)$, the natural
candidate detour source is the upstream junction of $e_1$ and the
candidate target is the downstream junction of $e_n$. However, if the
source junction has only one outgoing edge (the first closed edge
itself), no detour can diverge from it, the same problem arises
symmetrically at the target if it has only one incoming edge.

To handle this, Algorithm~\ref{alg:find-endpoints} walks upstream from the
source and downstream from the target until it reaches junctions with at least
two outgoing (resp.\ incoming) edges. Every intermediate edge encountered
during these walks is added to the blocked set, so the router cannot reuse it
as a shortcut.

\begin{algorithm}[H]
\caption{\textsc{FindDetourEndpoints}}
\label{alg:find-endpoints}
\begin{algorithmic}[1]
\Require closed edge sequence $E_c$, graph $G$
\Ensure source $s$, target $t$, extended blocked set $B$
\State start with $s$ and $t$ at the upstream and downstream ends of the closure
\State initialise $B$ with the closed edges
\While{$s$ has only one way out (the closure itself)}
  \State step one junction further upstream and add the traversed edge to $B$
\EndWhile
\While{$t$ has only one way in}
  \State step one junction further downstream and add the traversed edge to $B$
\EndWhile
\State \Return $s$, $t$, $B$
\end{algorithmic}
\end{algorithm}

Figure~\ref{fig:detour-endpoints} illustrates the procedure on a small graph
in which the natural endpoints $s_0$ and $t_0$ are both dead ends and the
walks promote $s$ and $t$ as the actual detour endpoints.

\begin{figure}[H]
\centering
\begin{tikzpicture}[
  >={Stealth[length=5pt]},
  node distance=1.5cm and 1.5cm,
  every node/.style={font=\small},
  jct/.style={circle, draw, fill=white, minimum size=8mm, inner sep=0pt, thick},
  endp/.style={circle, draw=ForestGreen!60!black, fill=ForestGreen!20,
               minimum size=8mm, inner sep=0pt, very thick},
  cand/.style={circle, draw=gray!70, dashed, fill=gray!8,
               minimum size=8mm, inner sep=0pt, thick},
  edge/.style={-{Stealth[length=5pt]}, thick, gray!55!black},
  closed/.style={-{Stealth[length=6pt]}, line width=1.4pt, red!75!black},
  blocked/.style={-{Stealth[length=5pt]}, thick, orange!85!black,
                  dash pattern=on 3pt off 2pt},
  walk/.style={-{Stealth[length=5pt]}, thick, blue!65!black, dashed},
  tag/.style={font=\scriptsize, inner sep=2pt}
]
  % main row of junctions (left to right)
  \node[jct]                  (A)  {$A$};
  \node[endp, right=of A]     (s)  {$s$};
  \node[cand, right=of s]     (s0) {$s_0$};
  \node[jct,  right=of s0]    (m)  {$M$};
  \node[cand, right=of m]     (t0) {$t_0$};
  \node[endp, right=of t0]    (t)  {$t$};
  \node[jct,  right=of t]     (Z)  {$Z$};

  % side branches that give s and t their second neighbour
  \node[jct, above=1.0cm of s] (X) {$X$};
  \node[jct, above=1.0cm of t] (Y) {$Y$};

  % free (untouched) edges
  \draw[edge] (A) -- (s);
  \draw[edge] (s) -- (X);
  \draw[edge] (Y) -- (t);
  \draw[edge] (t) -- (Z);

  % edges that are walked over and pushed into the blocked set
  \draw[blocked] (s)  -- (s0);
  \draw[blocked] (t0) -- (t);

  % closed edges (the closure itself)
  \draw[closed] (s0) -- (m);
  \draw[closed] (m)  -- (t0);

  % walk arcs above the row
  \draw[walk] (s0.north) to[out=135, in=45]
        node[above, tag] {walk upstream} (s.north);
  \draw[walk] (t0.north) to[out=135, in=45]
        node[above, tag] {walk downstream} (t.north);

  % role tags under each labelled node
  \node[tag, below=1pt of s,  text=ForestGreen!50!black] {final source};
  \node[tag, below=1pt of s0, text=gray!60!black]        {initial guess};
  \node[tag, below=1pt of t0, text=gray!60!black]        {initial guess};
  \node[tag, below=1pt of t,  text=ForestGreen!50!black] {final target};

  % legend
  \node[draw=gray!40, rounded corners=2pt, fill=gray!3,
        inner sep=4pt, font=\scriptsize, below=1.2cm of m, align=left] (leg) {%
    \begin{tabular}{@{}l@{\;\;}l@{}}
      \tikz[baseline=-0.6ex]\draw[closed]  (0,0) -- (0.6,0);  & closed edge ($\in E_c$) \\
      \tikz[baseline=-0.6ex]\draw[blocked] (0,0) -- (0.6,0);  & edge added to blocked set $B$ \\
      \tikz[baseline=-0.6ex]\draw[edge]    (0,0) -- (0.6,0);  & free edge \\
      \tikz[baseline=-0.6ex]\draw[walk]    (0,0) -- (0.6,0);  & walk performed by the algorithm \\
    \end{tabular}};
\end{tikzpicture}
\caption{Detour endpoint computation on an illustrative graph. The closure
$E_c = (s_0\!\to\!M,\ M\!\to\!t_0)$ is shown in red. Because $s_0$ has only
one outgoing edge (the closure itself) and $t_0$ has only one incoming edge,
neither can serve as a detour endpoint. The algorithm walks one junction
upstream from $s_0$ to $s$ (which has the additional outgoing edge
$s\!\to\!X$) and one junction downstream from $t_0$ to $t$ (which has the
additional incoming edge $Y\!\to\!t$). The two traversed edges (orange,
dashed) are appended to the blocked set $B$ so they cannot be reused by the
router as shortcuts around the closure.}
\label{fig:detour-endpoints}
\end{figure}

\subsubsection{Candidate generation}
\label{sssec:candidate_generation}

Two candidate-generation methods are implemented. Both start from the same
endpoint computation, but they use it differently. Yen's method keeps one
detour source $s$ and enumerates the $k$ best loopless paths from $s$ to
$t$, all of its plans have depth $d=0$. The upstream method starts
from the same effective source, then searches backwards up to
$d_{\max}$ hops. Each upstream node becomes a possible plan origin and the
shortest path from that origin to $t$ is retained if it exists. Yen's method
is the default in the end-to-end pipeline. The upstream method is used when
the experiment explicitly needs detours that begin before the closure point.

\paragraph{Yen's $k$-shortest loopless paths.}
\label{par:yen_dpm}

With source $s$, target $t$ and blocked set $B$ fixed, the first method
calls Yen's $k$-shortest loopless paths algorithm \cite{yen_algorithm} to
enumerate up to $k$ alternative routes from $s$ to $t$ that avoid every
edge in $B$.

As described in Algorithm~\ref{alg:yen},
Yen's algorithm extends Dijkstra's shortest-path search by iteratively
generating \emph{spur paths}. Let $P_1$ be the shortest path returned by
Dijkstra. At each subsequent iteration $i$, the algorithm considers every
prefix of the previously accepted path $P_{i-1}$ as a potential root
segment. For each root, it temporarily blocks the outgoing arc that was
used by any already-accepted path sharing that root, then runs Dijkstra
from the spur node (the last node of the root) to $t$. The resulting
spur path is concatenated with the root to form a complete candidate.
All candidates are stored in a min-heap keyed by total cost and the
minimum is promoted to $P_i$. A global signature set prevents the same
edge sequence from appearing twice. Heap entries carry a monotonically
increasing tie-breaking counter, so that two paths with identical cost
are ordered by insertion time rather than by a potentially fragile
lexicographic comparison of node lists.

\begin{algorithm}[H]
\caption{High-level pseudocode for \textsc{YenKShortestPaths}, which enumerates
up to $k$ loopless detour paths while avoiding blocked edges.}\label{alg:yen}
\begin{algorithmic}[1]
\Require graph $G$, source $s$, target $t$, integer $k \geq 1$, blocked set $B$
\Ensure list $\mathcal{P}$ of up to $k$ diverse shortest paths from $s$ to $t$ avoiding $B$
\State $P_1 \leftarrow$ shortest path from $s$ to $t$ avoiding $B$
\State $\mathcal{P} \leftarrow [P_1]$
\For{$i \leftarrow 2$ \textbf{to} $k$}
  \State \textbf{candidates} $\leftarrow \varnothing$
  \For{each prefix of the previously accepted path $P_{i-1}$}
    \State temporarily block the next arc taken by any already-accepted path that shares this prefix
    \State find the shortest path from the prefix's endpoint to $t$
    \State store \textit{prefix} $+$ \textit{new path} as a candidate
  \EndFor
  \State pick the cheapest unseen candidate as $P_i$ and append it to $\mathcal{P}$
  \State \textbf{stop} if no candidate remains
\EndFor
\State \Return $\mathcal{P}$
\end{algorithmic}
\end{algorithm}

\paragraph{Upstream search.}
\label{par:upstream}

The upstream method is based on the fact that vehicles approaching a closure
may arrive from several directions, not from one junction only. A reverse BFS
from $s$ records every reachable upstream node $u$ and its depth $d(u)$ up to
$d_{\max}$. The source $s$ itself is kept at depth $0$, so the method still
contains a fallback. For each discovered origin, Dijkstra is run
from that origin to $t$ while avoiding the blocked set $B$.

The returned plans are sorted by depth descending and then by total length.
This ordering is used later by the route rewriter: it tries the most upstream
applicable plan first, then falls back to shallower ones.
Algorithm~\ref{alg:upstream} summarizes this reverse search, path computation
and sorting step.

\begin{algorithm}[H]
\caption{High-level pseudocode for \textsc{UpstreamDeviationPlans}, which
searches upstream origins and sorts feasible detours by depth and length.}\label{alg:upstream}
\begin{algorithmic}[1]
\Require graph $G$, endpoints $s$, $t$, blocked set $B$, max depth $d_{\max}$
\Ensure deviation plans sorted by depth (desc), then length (asc)
\State reverse-BFS from $s$ up to $d_{\max}$ hops, tagging each reachable node $u$ with its depth $d(u)$
\State $\mathcal{P} \leftarrow [\,]$
\For{each upstream origin $u$ found (including $s$ at depth $0$)}
  \State compute the shortest path $\pi$ from $u$ to $t$ avoiding $B$
  \If{$\pi$ exists and is not a duplicate}
    \State save $\mathcal{D}(\pi,\, d(u))$ as a candidate plan in $\mathcal{P}$
  \EndIf
\EndFor
\State sort $\mathcal{P}$: deeper origins first, shorter paths first within each depth
\State \Return $\mathcal{P}$
\end{algorithmic}
\end{algorithm}

\noindent where $\mathcal{D}(\pi, d)$ denotes the deviation plan wrapping
path $\pi$ with BFS depth $d$.

Figure~\ref{fig:candidate-gen} illustrates the two methods on the same toy
network: nodes $s$ and $t$ are the detour endpoints, $p$ and $q$ are
intermediate junctions and the direct edge $s\to t$ belongs to the blocked
set $B$.

\begin{figure}[H]
\centering
%%------------------------------------------------------------
%% (a) Yen's k-shortest paths
%%------------------------------------------------------------
\begin{minipage}[t]{0.44\linewidth}
\centering
\begin{tikzpicture}[
  junc/.style={circle, draw=gray!60, fill=gray!12, thick,
               minimum size=8mm, inner sep=0pt, font=\small\sffamily},
  endpt/.style={junc, draw=blue!55!black, fill=blue!12},
  cl/.style={-{Stealth[length=6pt]}, red!68, dashed, line width=1.2pt},
  base/.style={-{Stealth[length=5pt]}, gray!42, thin},
  pa/.style={-{Stealth[length=6pt]}, blue!65, line width=1.6pt},
  pb/.style={-{Stealth[length=6pt]}, orange!85!black, line width=1.6pt},
]
  \node[endpt] (s) at (0,0)     {$s$};
  \node[junc]  (p) at (2.0,1.3) {$p$};
  \node[junc]  (q) at (2.0,-1.3){$q$};
  \node[endpt] (t) at (4.0,0)   {$t$};

  %% base network
  \draw[base] (s) -- (p);\quad\draw[base] (p) -- (t);
  \draw[base] (s) -- (q);\quad\draw[base] (q) -- (t);

  %% blocked direct edge
  \draw[cl] (s) -- node[above, font=\scriptsize, red!68]{blocked} (t);

  %% P1: s -> p -> t
  \draw[pa] (s) to[bend left=18] (p);
  \draw[pa] (p) to[bend left=18] (t);

  %% P2: s -> q -> t
  \draw[pb] (s) to[bend right=18] (q);
  \draw[pb] (q) to[bend right=18] (t);

  %% inline legend
  \draw[pa, line width=1.2pt] (-0.1,-2.1)--(0.6,-2.1)
    node[right, font=\scriptsize]{$P_1\colon s\!\to\!p\!\to\!t$};
  \draw[pb, line width=1.2pt] (-0.1,-2.55)--(0.6,-2.55)
    node[right, font=\scriptsize]{$P_2\colon s\!\to\!q\!\to\!t$};
  \draw[cl,  line width=1.2pt] (-0.1,-3.0) --(0.6,-3.0)
    node[right, font=\scriptsize]{blocked ($B$)};
\end{tikzpicture}

\smallskip\small (a)~Yen's $k$-shortest paths ($k=2$)
\end{minipage}%
\hfill
%%------------------------------------------------------------
%% (b) Upstream multi-origin
%%------------------------------------------------------------
\begin{minipage}[t]{0.52\linewidth}
\centering
\begin{tikzpicture}[
  junc/.style={circle, draw=gray!60, fill=gray!12, thick,
               minimum size=8mm, inner sep=0pt, font=\small\sffamily},
  endpt/.style={junc, draw=blue!55!black, fill=blue!12},
  upst/.style={junc, draw=teal!60!black, fill=teal!10},
  cl/.style={-{Stealth[length=6pt]}, red!68, dashed, line width=1.2pt},
  base/.style={-{Stealth[length=5pt]}, gray!42, thin},
  pa/.style={-{Stealth[length=6pt]}, blue!65, line width=1.6pt},
  pb/.style={-{Stealth[length=6pt]}, orange!85!black, line width=1.6pt},
]
  \node[endpt] (s)  at (0,0)      {$s$};
  \node[junc]  (p)  at (2.0,1.3)  {$p$};
  \node[junc]  (q)  at (2.0,-1.3) {$q$};
  \node[endpt] (t)  at (4.0,0)    {$t$};
  \node[upst]  (u1) at (-1.8,0)   {$u_1$};
  \node[upst]  (u2) at (-1.8,-1.3){$u_2$};

  %% base network
  \draw[base] (s)  -- (p);\quad\draw[base] (p) -- (t);
  \draw[base] (s)  -- (q);\quad\draw[base] (q) -- (t);
  \draw[base] (u1) -- (s);   %% u1 -> s
  \draw[base] (u2) -- (s);   %% u2 -> s (exists, not used by u2's path)
  \draw[base] (u2) to[bend right=12] (q); %% shortcut u2 -> q

  %% blocked direct edge
  \draw[cl] (s) -- node[above, font=\scriptsize, red!68]{blocked} (t);

  %% BFS boundary
  \begin{scope}[on background layer]
    \draw[teal!45, dashed, rounded corners=5pt, line width=0.8pt]
      (-2.45,-1.9) rectangle (-1.15,0.55);
  \end{scope}
  \node[font=\scriptsize\itshape, teal!60!black] at (-1.8,0.9)
    {upstream$(s,\,1)$};

  %% path from u1: u1 -> s -> p -> t
  \draw[pa] (u1) -- (s);
  \draw[pa] (s)  to[bend left=18]  (p);
  \draw[pa] (p)  to[bend left=18]  (t);

  %% path from u2: u2 -> q -> t  (bypasses s)
  \draw[pb] (u2) to[bend right=20] (q);
  \draw[pb] (q)  to[bend right=18] (t);

  %% inline legend
  \draw[pa, line width=1.2pt] (-0.1,-2.1)--(0.6,-2.1)
    node[right, font=\scriptsize]{$u_1\colon u_1\!\to\!s\!\to\!p\!\to\!t$};
  \draw[pb, line width=1.2pt] (-0.1,-2.55)--(0.6,-2.55)
    node[right, font=\scriptsize]{$u_2\colon u_2\!\to\!q\!\to\!t$ \emph{(bypasses $s$)}};
  \draw[cl,  line width=1.2pt] (-0.1,-3.0) --(0.6,-3.0)
    node[right, font=\scriptsize]{blocked ($B$)};
\end{tikzpicture}

\smallskip\small (b)~Upstream multi-origin ($d=1$)
\end{minipage}

\caption{Comparison of the two candidate generation methods on a toy
network. In both cases the edge $s\!\to\!t$ is blocked (red dashed).
\textit{(a)}~Yen's algorithm produces $k=2$ loopless paths, both originating
from $s$: $P_1$ takes the upper corridor through $p$ and $P_2$ the lower
corridor through $q$.
\textit{(b)}~The upstream BFS at depth $d=1$ identifies two origins
$\{u_1, u_2\}$ (teal). $u_1$ routes through $s$ and then takes the upper
corridor, while $u_2$ holds a direct connection to $q$, so its optimal
path bypasses $s$ entirely, a candidate that Yen's method, being anchored
at $s$, would never generate.}
\label{fig:candidate-gen}
\end{figure}

Each accepted path is wrapped in a \texttt{DeviationPlan} object. The
object records the plan identifier, its rank within the $k$-shortest
list, the ordered sequences of junction and edge identifiers and the
total geometric length in meters (computed by summing the
\texttt{length\_m} attribute of every edge on the path).

No post-filtering is applied at this point. Plans are not removed because of a
maximum detour-length ratio and near-duplicates are not removed by edge-set
overlap. This keeps the planning stage simple and avoids discarding a route
before it has been tested in simulation. Two routes can share most of their
edges and still behave differently if they diverge near a bottleneck or pass
through a different sequence of signalised intersections. The Monte Carlo
evaluation is used later to separate useful candidates from poor ones.

\subsubsection{Route rewriting}

Detours are applied by rewriting route files before SUMO starts, rather than
by using SUMO's dynamic
\emph{rerouters} \cite{sumo_rerouter}.
This keeps the experiment static and inspectable: \texttt{duarouter} already
provides route-choice variation through randomized edge weights and the final
route file can be checked before the simulation.

After each Monte Carlo call to \texttt{duarouter}, the rewriter creates one
\path|routes_{plan_id}.xml| file per candidate plan. Only vehicles that
traverse a closed edge are modified. For each affected vehicle, candidate plans
are tried from deepest to shallowest and the first plan whose origin appears
upstream of the closure in that vehicle's route is selected
(Algorithm~\ref{alg:rewrite}). Yen plans have only the evaluated depth-0
candidate, upstream plans also carry the depth-0 fallback.

The selected plan $\mathcal{D}^*$ is inserted by replacing the route segment
from $\mathcal{D}^*.s$ to the last closed edge with the detour edge sequence
$\mathcal{D}^*.P$. This is a splice in the vehicle's edge list.

\begin{algorithm}[H]
\caption{High-level pseudocode for \textsc{DepthAwareRewrite}, which selects
the deepest applicable detour and splices it into an affected route.}\label{alg:rewrite}
\begin{algorithmic}[1]
\Require vehicle route $R$, candidate plans $\mathcal{D}$ sorted by depth (desc),
         closed edge set $B_c$
\Ensure rewritten route $R'$
\If{$R$ does not traverse any edge in $B_c$}
  \State \Return $R$\hfill\textit{// route unaffected by the closure}
\EndIf
\State pick the deepest plan $\mathcal{D}^*$ whose origin $s$ lies on $R$
       upstream of the closure
       \hfill\textit{// shallower plans provide deterministic fallback}
\State splice $\mathcal{D}^*.P$ into $R$, replacing the segment that runs
       from the edge departing $\mathcal{D}^*.s$ up to the last closed edge
\State \Return the spliced route
\end{algorithmic}
\end{algorithm}

The baseline route file is left unchanged, each candidate is evaluated through
its own rewritten copy.

Figure~\ref{fig:route-rewrite} illustrates this route-list splice for one
affected vehicle.

\begin{figure}[H]
  \centering

  \begin{tikzpicture}[
    ebox/.style={
      draw,
      rounded corners=3pt,
      fill=gray!10,
      thick,
      minimum width=1.05cm,
      minimum height=0.65cm,
      align=center,
      font=\scriptsize\ttfamily
    },
    cbox/.style={
      draw=red!65!black,
      fill=red!12,
      rounded corners=3pt,
      thick,
      minimum width=1.05cm,
      minimum height=0.65cm,
      align=center,
      font=\scriptsize\ttfamily
    },
    dbox/.style={
      draw=teal!65!black,
      fill=teal!12,
      rounded corners=3pt,
      thick,
      minimum width=1.05cm,
      minimum height=0.65cm,
      align=center,
      font=\scriptsize\ttfamily
    },
    arr/.style={-{Stealth[length=5pt]}, thick},
    node distance=0.22cm
  ]

    % ---- Original route (top row) ----
    \node[ebox] (e1) {$e_1$};
    \node[ebox, right=of e1] (e2) {$e_2$};
    \node[ebox, right=of e2] (e3) {$e_3$};
    \node[cbox, right=of e3] (ec) {$e_c$};
    \node[ebox, right=of ec] (e4) {$e_4$};
    \node[ebox, right=of e4] (e5) {$e_5$};

    \node[
      left=0.55cm of e1,
      font=\small,
      align=right,
      text width=3cm
    ] {Original route:};

    \draw[arr] (e1) -- (e2);
    \draw[arr] (e2) -- (e3);
    \draw[arr, red!50] (e3) -- (ec);
    \draw[arr, red!50] (ec) -- (e4);
    \draw[arr] (e4) -- (e5);

    % Closed-edge annotation
    \draw[
      red!60,
      decorate,
      decoration={brace, amplitude=5pt, raise=3pt}
    ]
      (ec.north west) -- (ec.north east)
      node[midway, above=8pt, font=\scriptsize, red!75!black]{closed};

    % ---- Central arrow ----
    \node[font=\small, align=center] (lbl) at (5.8, -1.1)
      {\textsc{ApplyDetour}($\mathcal{D}^*$)};

    \draw[arr, blue!55, thick] (3.8, -0.42) -- (3.8, -1.75);

    % ---- Rewritten route (bottom row) ----
    \node[ebox, below=2cm of e1] (f1) {$e_1$};
    \node[ebox, right=of f1] (f2) {$e_2$};
    \node[ebox, right=of f2] (f3) {$e_3$};
    \node[dbox, right=of f3] (d1) {$d_1$};
    \node[dbox, right=of d1] (d2) {$d_2$};
    \node[ebox, right=of d2] (f4) {$e_4$};
    \node[ebox, right=of f4] (f5) {$e_5$};

    \node[
      left=0.55cm of f1,
      font=\small,
      align=right,
      text width=3cm
    ] {Rewritten route:};

    \draw[arr] (f1) -- (f2);
    \draw[arr] (f2) -- (f3);
    \draw[arr, teal!65!black] (f3) -- (d1);
    \draw[arr, teal!65!black] (d1) -- (d2);
    \draw[arr, teal!65!black] (d2) -- (f4);
    \draw[arr] (f4) -- (f5);

    % Detour annotation
    \draw[
      teal!65!black,
      decorate,
      decoration={brace, amplitude=5pt, raise=3pt}
    ]
      (d1.north west) -- (d2.north east)
      node[midway, above=8pt, font=\scriptsize, teal!70!black]
      {detour $\mathcal{D}^*\!$};

  \end{tikzpicture}

  \caption{Route rewriting for one affected vehicle. The closed edge $e_c$
  (red) is replaced by the selected detour sequence $(d_1, d_2)$ (teal), while
  the unaffected prefix and suffix are preserved.}

  \label{fig:route-rewrite}
\end{figure}
